\chapter{الأعداد حتى 10\,000}\label{ch:g3:numbers}

ألف، وألفان، وثلاثة آلاف\,\dots\ الأعداد الكبيرة تُبنى من الصغيرة،
بفكرة بديعة: \emph{مرتبة} الرقم هي التي تحدد قيمته. يعلّم هذا
الفصل كيف نقرأ الأعداد حتى $10\,000$ ونكتبها ونقارنها
ونرتّبها.

\section{قراءة الأعداد وكتابتها}

\begin{definition}[القيمة المكانية]\label{def:g3:numbers:place}
يُكتب العدد بالأرقام $0, 1, 2, \dots, 9$. وبالقراءة من
اليمين تكون المراتب: آحاد، وعشرات، ومئات، وآلاف. ففي
$3\,254$:
\[
3\,254 = 3 \text{ آلاف} + 2 \text{ مئات} + 5 \text{ عشرات}
+ 4 \text{ آحاد}.
\]
\end{definition}

\begin{omfigure}
\begin{tikzpicture}[scale=0.9]
  \draw (0,0) grid[xstep=2.3, ystep=0.85] (9.2,1.7);
  \node[font=\small] at (1.15,1.275) {آلاف};
  \node[font=\small] at (3.45,1.275) {مئات};
  \node[font=\small] at (5.75,1.275) {عشرات};
  \node[font=\small] at (8.05,1.275) {آحاد};
  \node[font=\large, omDef] at (1.15,0.425) {$3$};
  \node[font=\large, omDef] at (3.45,0.425) {$2$};
  \node[font=\large, omDef] at (5.75,0.425) {$5$};
  \node[font=\large, omDef] at (8.05,0.425) {$4$};
\end{tikzpicture}

{\small جدول المراتب للعدد $3\,254$، ويُقرأ «ثلاثة آلاف
ومئتان وأربعة وخمسون».}
\end{omfigure}

\begin{example}[الصفر الذي يحفظ المرتبة]\label{ex:g3:numbers:zero}
«أربعة آلاف وسبعة» يُكتب $4\,007$: الرقم $4$ في
الآلاف، والرقم $7$ في الآحاد، و\emph{أصفار} بينهما تُبقي
كل رقم في مرتبته الصحيحة. ولولاها لكان $47$ عددًا
مختلفًا جدًّا!
\end{example}

\begin{example}[التبديل]\label{ex:g3:numbers:exchange}
عشرة آحاد تصنع عشرة؛ وعشر عشرات تصنع مئة؛ وعشر مئات تصنع
ألفًا. إذن $10$ مئات $= 1\,000$، و $32$ عشرة
$= 320$: فالعدد $320$ يحتوي $32$ عشرة (لا الرقم
$2$ وحده!).
\end{example}

\section{المقارنة والترتيب}

\begin{method}[مقارنة عددين]\label{met:g3:numbers:compare}
\begin{enumerate}
  \item عُدّ الأرقام: فالأرقام الأكثر تعني عددًا أكبر
        ($1\,002 > 998$).
  \item وعند تساوي عدد الأرقام: قارن رقمًا رقمًا بدءًا من
        \emph{اليسار}. وأول اختلاف هو الذي يقرر:
        $5\,614 > 5\,589$ لأن $6 > 5$ في المئات.
  \item اُكتب الجواب بالإشارتين $<$ (أصغر من) أو $>$
        (أكبر من): فالجهة المفتوحة تواجه دائمًا العدد الأكبر.
\end{enumerate}
\end{method}

\begin{omfigure}
\begin{tikzpicture}[scale=1.1]
  \draw[-Stealth] (-0.3,0) -- (10.5,0);
  \foreach \i in {0,1,...,10} \draw (\i,-0.08) -- (\i,0.08);
  \node[below, font=\small] at (0,-0.1) {$0$};
  \node[below, font=\small] at (5,-0.1) {$500$};
  \node[below, font=\small] at (10,-0.1) {$1000$};
  \fill[omThm] (2.4,0) circle (2pt) node[above, font=\small] {$240$};
  \fill[omDef] (6.9,0) circle (2pt) node[above, font=\small] {$690$};
  \fill[omMeth] (9.05,0) circle (2pt) node[above, font=\small] {$905$};
\end{tikzpicture}

{\small أعداد على مستقيم مدرَّج بالمئات: كلما اتجهت يمينًا كبر
العدد. $240 < 690 < 905$.}
\end{omfigure}

\begin{example}\label{ex:g3:numbers:order}
رتّب من الأصغر إلى الأكبر: $780$ و $8\,700$ و $807$ و $87$.
أولًا بعدد الأرقام: $87$ (رقمان)، ثم $780$ و $807$
(ثلاثة)، ثم $8\,700$ (أربعة). وبين $780$ و $807$: قارن رقمَي
المئات، $7 < 8$، إذن $780 < 807$. والترتيب النهائي:
\[
87 < 780 < 807 < 8\,700 .
\]
\end{example}

\section{الزوجي والفردي}

\begin{definition}[العدد الزوجي والعدد الفردي]\label{def:g3:numbers:evenodd}
نقول عن العدد إنه \emph{عدد زوجي}\index{عدد زوجي} إذا أمكن
تقسيمه إلى جزأين صحيحين متساويين --- ورقم آحاده $0$ أو $2$ أو $4$
أو $6$ أو $8$. وإلا فهو \emph{عدد فردي}\index{عدد فردي} (ورقم
آحاده $1$ أو $3$ أو $5$ أو $7$ أو $9$).
\end{definition}

\begin{example}
العدد $46$ زوجي ($46 = 23 + 23$)؛ والعدد $47$ فردي (فالجزآن
المتساويان يحتاجان إلى نصف شيء). ولا يهمّ إلا الرقم
\emph{الأخير}: فالعدد $3\,578$ زوجي، والعدد $2\,341$ فردي.
\end{example}

\begin{omfigure}
\begin{tikzpicture}[scale=0.45]
  \foreach \i in {0,...,5} \fill[omDef] ({\i},0) circle (0.28);
  \foreach \i in {0,...,5} \fill[omDef] ({\i},1) circle (0.28);
  \node[below, font=\small] at (2.5,-0.7) {$12$: زوجي، صفّان متساويان};
  \begin{scope}[xshift=10.5cm]
  \foreach \i in {0,...,4} \fill[omThm] ({\i},0) circle (0.28);
  \foreach \i in {0,...,4} \fill[omThm] ({\i},1) circle (0.28);
  \fill[omThm] (5,0) circle (0.28);
  \draw[omExo, thick] (5,0) circle (0.5);
  \node[below, font=\small] at (2.5,-0.7) {$11$: فردي، نقطة وحيدة};
  \end{scope}
\end{tikzpicture}

{\small \omterm{def:g3:numbers:evenodd}{الأعداد الزوجية} تنقسم إلى صفّين متساويين؛ أما \omterm{def:g3:numbers:evenodd}{الأعداد الفردية}
فتترك واحدة وحيدة دائمًا.}
\end{omfigure}

\section{تمارين}

\begin{exercise}[$\star$]\label{exo:g3:numbers:1}
اُكتب بالأرقام: «ألفان وثلاثمئة وخمسة وأربعون»؛ «ستة
آلاف وخمسون»؛ «تسعة آلاف وواحد».
\end{exercise}

\begin{exercise}[$\star$]\label{exo:g3:numbers:2}
اُكتب بالحروف: $1\,208$؛ \quad $4\,070$؛ \quad $9\,999$.
\end{exercise}

\begin{exercise}[$\star$]\label{exo:g3:numbers:3}
في $7\,382$: أي رقم في مرتبة العشرات؟ وأيّها في مرتبة
الآلاف؟ وكم \emph{مئة} يحتوي العدد في المجموع (واِنتبه: ليس رقم
المئات وحده --- وفكّر في
\cref{ex:g3:numbers:exchange})؟
\end{exercise}

\begin{exercise}[$\star$]\label{exo:g3:numbers:4}
فكّك كما في \cref{def:g3:numbers:place}:
$5\,631$؛ \quad $2\,046$؛ \quad $8\,500$.
\end{exercise}

\begin{exercise}[$\star$]\label{exo:g3:numbers:5}
اُنقل وأكمل بالإشارة $<$ أو $>$:
\[
456 \;?\; 465, \qquad
1\,203 \;?\; 989, \qquad
6\,090 \;?\; 6\,009, \qquad
9\,999 \;?\; 10\,000 .
\]
\end{exercise}

\begin{exercise}[$\star$]\label{exo:g3:numbers:6}
رتّب من الأصغر إلى الأكبر: $530$؛ \ $3\,500$؛ \ $353$؛ \ $5\,033$؛
\ $503$.
\end{exercise}

\begin{exercise}[$\star$]\label{exo:g3:numbers:7}
على مستقيم مدرَّج بالمئات من $0$ إلى $1\,000$، ضع
(تقريبًا) الأعداد $150$ و $480$ و $820$ و $990$.
\end{exercise}

\begin{exercise}[$\star$]\label{exo:g3:numbers:8}
زوجي أم فردي؟ \ $34$؛ \ $57$؛ \ $70$؛ \ $2\,481$؛ \ $6\,548$.
\end{exercise}

\begin{exercise}[$\star$]\label{exo:g3:numbers:9}
عُدّ بالخطوات: من $2\,970$، عُدّ خمس خطوات من $10$؛ ومن $860$،
عُدّ أربع خطوات من $100$؛ ومن $9\,996$، عُدّ أربع خطوات من $1$.
فماذا تلاحظ في نهاية الأخيرة؟
\end{exercise}

\begin{exercise}[$\star\star$]\label{exo:g3:numbers:10}
مستعملًا كلًّا من الأرقام $3$ و $7$ و $0$ و $5$ مرة واحدة
بالضبط، اُكتب أكبر عدد ممكن من أربعة أرقام، ثم أصغره (ولا يبدأ
العدد بالرقم $0$).
\end{exercise}

\begin{exercise}[$\star\star$]\label{exo:g3:numbers:11}
أنا عدد من ثلاثة أرقام. أرقامي هي $2$ و $5$ و $8$، وأنا زوجي،
وأنا أكبر من $800$. فمن أنا؟
\end{exercise}
